\documentclass[11pt,letterpaper]{article}
\PassOptionsToPackage{hyphens}{url}
\usepackage[margin=.83in,top=.8in,bottom=.8in,headheight=15pt]{geometry}
\usepackage[T1]{fontenc}
\usepackage{lmodern,amsmath,amssymb,mathtools,bm,booktabs,array,graphicx,tikz}
\usetikzlibrary{arrows.meta,positioning,calc,decorations.pathreplacing,patterns}
\definecolor{ink}{HTML}{17344D}\definecolor{accent}{HTML}{087E8B}\definecolor{pale}{HTML}{EDF5F5}\definecolor{gold}{HTML}{AD6B19}
\providecommand{\dd}{\,\mathrm{d}}\providecommand{\R}{\mathbb{R}}
\providecommand{\E}{\mathbb{E}}\providecommand{\Var}{\operatorname{Var}}
\providecommand{\Cov}{\operatorname{Cov}}\providecommand{\diag}{\operatorname{diag}}
\providecommand{\norm}[1]{\left\lVert#1\right\rVert}
\providecommand{\abs}[1]{\left|#1\right|}
\usepackage{microtype,hyperref,fancyhdr,enumitem,amsthm,sectsty,needspace}
\allsectionsfont{\raggedright}
\hypersetup{colorlinks=true,linkcolor=ink,urlcolor=accent,pdfauthor={Zan Ahmad},pdftitle={Finite-population genetics: inheritance, drift, and fixation / Review}}
\setlength{\parindent}{0pt}\setlength{\parskip}{6pt}\linespread{1.07}
\setlist{itemsep=3pt,topsep=4pt,leftmargin=*}
\emergencystretch=2em\clubpenalty=10000\widowpenalty=10000
\pagestyle{fancy}\fancyhf{}\fancyhead[L]{\small Mathematical Modeling}\fancyhead[R]{\small Lecture 7}
\fancyfoot[L]{\footnotesize Zan Ahmad}\fancyfoot[R]{\footnotesize\thepage}\renewcommand{\headrulewidth}{.3pt}
\newtheorem{proposition}{Proposition}\newtheorem{theorem}{Theorem}\newtheorem{lemma}{Lemma}
\theoremstyle{definition}\newtheorem{definition}{Definition}\newtheorem{example}{Example}
\numberwithin{equation}{section}
\begin{document}
\section*{Finite-population genetics: inheritance, drift, and fixation --- review sheet}
Use this after attempting the worked examples and practice. Explain each claim in words before calculating.

\section*{Key ideas}
\begin{itemize}
\item Diploid allele frequencies count copies; genotype frequencies count individuals.
\item Independent gametes give probabilities p$^2$, 2p(1-p), (1-p)$^2$.
\item Neutral drift can change one trajectory while preserving ensemble expected frequency.
\item At fixed p, one-generation frequency variance scales as 1/(2N).
\item Without mutation or migration, zero and one are absorbing boundaries.
\end{itemize}
\section*{Common errors and repairs}
\textbf{Using AB frequency as A allele frequency.} Count two A copies per AA and one per AB, then divide by 2N.

\textbf{Interpreting the dashed expectation as a required trajectory.} One finite-population realization can differ; even a finite ensemble mean fluctuates.

\textbf{Declaring every upward fluctuation evidence of selection.} Neutral random sampling can produce upward changes; compare competing explanations with appropriate data.

\section*{Retrieve without your notes}
\begin{itemize}
\item At p =\allowbreak{} 0.5, what are the three genotype probabilities?
\item At fixed p, what happens to one-generation variance when N doubles?
\item Why does sampling without replacement from the entire pool fail to create the modeled drift?
\end{itemize}
\section*{Readiness check}
\begin{itemize}
\item Compute allele frequency from diploid genotype counts without confusing copies and individuals.
\item Derive the three offspring genotype probabilities for independent gametes.
\item Explain the difference between an expected frequency and one finite-population trajectory.
\item Predict how population size changes one-generation sampling variation.
\item Check probability totals, attainable frequencies and absorbing limits.
\end{itemize}
\end{document}
